\section{Introduction}
\label{intro}
We consider nonparametric density estimation on a compact support. Given a random variable $X \sim P_0$, where $P_0$ is absolutely continuous, and i.i.d. samples $x_{1:n}$ from $P_0$, our goal is to estimate the true density function $p_0$ of the underlying distribution $P_0$. This paper aims to provide a thorough theoretical analysis of univariate density estimation with a variational penalty and its application. The statistical model for $P_0$ will be nonparametric up till assuming that $p_0$ has support on a known interval $[a,b]$ and is a cadlag function with a bounded variation norm, explained in detail below. Our framework is naturally extended to the multivariate case. 

Nonparametric density estimation methods typically include kernel-based approaches, splines, and wavelet techniques. Kernel density estimation (KDE), despite its simplicity, suffers from several drawbacks: it poorly captures densities with rapidly varying regions and faces severe challenges due to the curse of dimensionality in the multivariate case.

The first drawback of KDE arises from the fact that it is a linear smoother whose fitted values depend linearly on the observed responses. TV penalties have been introduced
into spline regression to resolve the same problem with kernel regression, smoothing splines, etc.    \citet{mammen1997locally} developed local adaptive splines (LAS) demonstrating $L^2$ convergence, while \citet{tibshirani2014adaptive} formulated restricted LAS and trend filtering (TF) as general Lasso regressions. Later, the TV penalty was also introduced into the density estimation problem \citep{bak2021penalized, sadhanala2024exponential}. The TV-penalized logspline density estimation (PLSDE) proposed in \citet{bak2021penalized} is more relevant to our problem, and it intends to tackle the oscillation issue of logspines method \citep{kooperberg1991study, kooperberg1992logspline}. \citet{kooperberg1992logspline} employs an exponential link transformation ensuring positivity:
\begin{align}
\label{Link_function}
p_{f}(x) = \frac{\exp(f(x))}{\int_a^b \exp(f(u)) d\mu(u)},
\end{align}
with splines $f(x) = \sum_j \beta_j \phi_j(x)$, typically cubic B-splines. Parameters are estimated via maximum likelihood, often guided by criteria such as AIC or BIC. A theoretical analysis of the logspline model is provided in \citet{stone1990large}.
\citet{bak2021penalized} employed TV penalty and BIC criteria to provide univariate KL-divergence convergence analysis and its generalization to bivariate case. However, their method, Penalized Log-spline Density Estimation (PLSDE), encounters difficulties generalizing to multivariate settings without assuming higher-order continuity. This is a manifestation of the curse of dimensionality, and it falls into the same problem for the multivariate spline when the TV penalty is not introduced. Another problem is that PLSDE is limited to the uniform knots, which is not preferred in the multivariate case. An argument of the knot placement problem is provided in Section~\ref{The Highly Adaptive Lasso (HAL) Assumptions}.

In contrast, the Highly Adaptive Lasso (HAL) proposed by \citet{van2017generally} and practically introduced in  \citet{benkeser2016highly} estimates multivariate càdlàg functions defined on $[0,1]^d$ using a sectional variation norm, yielding uniform consistency and pointwise asymptotic normality without dimension-enforced smoothness assumptions. The same function class has received broad attention in other settings to avoid the curse of dimensionality asymptotically, including nonparametric isotonic/variation-constrained regression \citep{fang2021multivariate} and MARS-type regression \citep{ki2024mars}; in this paper, we focus on nonparametric density estimation. The density estimator based on HAL is referred to as HAL-MLE log-splines method. In this paper, we restrict our focus to univariate HAL-MLE in order to compare the performance with the classical log-spines methods and demonstrate theoretical connections to LAS and TF approaches, even though the HAL theory applies to the multivariate case. Other than density estimation, HAL-MLE also provide the asymptotic efficiency guarantee for general pathwise-differentiable statistical estimand, i.e. moments, survival probability, or percentiles, by a simple plug-in or a single-step TMLE procedure as developed in \citet{van2017generally, van2023efficient}.

The remainder of this paper is organized as follows. Section~\ref{The Highly Adaptive Lasso (HAL) Assumptions} introduces the HAL assumption and establishes its connection to the classical bounded total variation (BTV) assumption underlying local adaptive splines. Section~\ref{HAL Density Estimation} presents the construction of HAL-MLE with the log-spline link function \citep{leonard1978density, silverman1982estimation, kooperberg1992logspline, rytgaard2023estimation}. Section~\ref{Theoretical Properties} presents the theoretical results on its univariate $L^2$ convergence, asymptotic linearity, pointwise asymptotic normality, and uniform convergence. We also propose a variance estimator for the density based on the delta method. Section~\ref{section:HAL-MLE for Pathwise Differentiable Statistical Estimands} considers the plug-in HAL-MLE and HAL-TMLE of pathwise differentiable statistical estimands, shown to achieve asymptotic efficiency with influence-curve-based variance estimation. Section~\ref{sec:simulation} reports simulation results, empirically verifying the theoretical guarantees mentioned in the previous sections, and comparing the finite sample performance of HAL-MLE with TF (applied to density estimation), log-splines, and KDE. In Section~\ref{sec:case-study}, we present a case study with the galaxy velocity data to visualize the convergence, confidence intervals, and targetings. Additional computational details and optimization algorithm performance analyses are provided in Section S.10 of the Supplementary Material \citep{hou2026supplement}. We provide a Python package for HAL-MLE density estimation, available at \url{https://github.com/zhengpu-berkeley/HALDensity}, and experiment details at \url{https://github.com/yilongHou/link_HAL_MLE}.
